\documentclass[11pt]{article}

\usepackage[margin=1in]{geometry}
\usepackage{amsmath}
\usepackage{amssymb}
\usepackage{enumitem}
\usepackage[hidelinks]{hyperref}

\setlist[itemize]{leftmargin=*, itemsep=0.5em}
\setlist[enumerate]{leftmargin=*, itemsep=0.5em}

\title{Smoothed Particle Hydrodynamics Study Notes}
\author{}
\date{}

\begin{document}

\maketitle
\tableofcontents

\section{Fluid Simulation Methods and Where SPH Fits}

In numerical fluid simulation, continuous fields such as velocity, density,
pressure, and temperature must be represented by finitely many degrees of
freedom. Methods are often classified by how those degrees of freedom represent
space and how they move relative to the fluid. The broadest distinction is
between \textbf{Eulerian} methods, which observe fluid through fixed locations,
and \textbf{Lagrangian} methods, which follow moving parcels of material.
Hybrid methods combine both viewpoints.

\subsection{Eulerian Mesh-Based Methods}

Eulerian methods divide the spatial domain into a structured or unstructured
mesh and store physical quantities at grid nodes, cell centers, or element
faces. Fluid flows through this fixed spatial representation. Common examples
include:

\begin{itemize}
    \item \textbf{Finite Difference Methods (FDM):} Approximate derivatives
    directly from values on a usually structured grid. They are conceptually
    simple and efficient, but are less convenient for complicated geometries.
    \item \textbf{Finite Volume Methods (FVM):} Integrate the governing
    equations over control volumes and compute fluxes through cell faces.
    Their local conservation properties make them especially common in
    computational fluid dynamics (CFD).
    \item \textbf{Finite Element Methods (FEM):} Represent the solution with
    basis functions over mesh elements. They naturally support irregular
    geometries, variable resolution, and rigorous weak formulations.
    \item \textbf{Lattice Boltzmann Methods (LBM):} Evolve discrete particle
    distribution functions on a lattice rather than discretizing the
    Navier--Stokes equations directly. They are highly parallel and useful for
    complex or porous boundaries, although standard formulations use regular
    lattices and have restrictions on compressibility and stability.
\end{itemize}

Mesh-based methods generally provide accurate differential operators, mature
pressure solvers, and efficient computation on regular grids. They are often
preferred for flows in fixed domains and for high-accuracy engineering CFD.
Their main difficulty is that the fluid must be advected through the grid,
which can introduce numerical diffusion. Free surfaces require an additional
interface representation, such as a level-set or volume-of-fluid method. A
fixed Eulerian mesh itself does not deform, but moving boundaries and interfaces
can still require cut-cell, immersed-boundary, adaptive-mesh, or remeshing
techniques.

\subsection{Lagrangian and Mesh-Free Methods}

Lagrangian methods move their computational degrees of freedom with the
material. \textbf{Smoothed Particle Hydrodynamics (SPH)} is both a Lagrangian
particle method and a mesh-free method: it represents the fluid by particles
without permanent connectivity between them. Each particle carries quantities
such as mass, position, velocity, density, and pressure. Field values and their
derivatives are estimated from weighted sums over neighboring particles using
a smoothing kernel.

Other methods in or near this category include \textbf{vortex methods}, which
represent vorticity using moving particles or filaments, and other mesh-free
collocation or moving-least-squares methods. These methods differ from SPH in
what their particles represent and how they reconstruct continuous fields.

\subsection{Hybrid Particle--Grid Methods}

Particle-in-cell methods use particles to transport material information and a
background grid to compute forces or solve pressure. Important variants include
\textbf{PIC}, \textbf{FLIP}, and \textbf{APIC}. The \textbf{Material Point
Method (MPM)} uses a related particle--grid transfer for fluids, solids, and
coupled materials. These approaches seek a compromise: particles handle
advection and changing topology, while the grid supplies structured and often
more stable differential operators. Their disadvantages include transfer error
between particles and the grid, grid-crossing artifacts, and the complexity of
maintaining two representations.

\subsection{Benefits of SPH}

\begin{itemize}
    \item \textbf{Natural free surfaces:} A fluid boundary is simply the edge
    of the particle distribution, so no separate interface-tracking field is
    required for many applications.
    \item \textbf{Large deformation and topology changes:} Splashing,
    fragmentation, merging, and moving boundaries do not distort a mesh or
    require connectivity to be rebuilt.
    \item \textbf{Lagrangian advection:} Material quantities move with the
    particles, reducing the numerical diffusion associated with transporting
    them through a fixed grid.
    \item \textbf{Conservation and coupling:} Symmetric pairwise force
    formulations can conserve mass exactly and conserve linear and angular
    momentum well. The particle representation also couples naturally to
    moving rigid bodies and deformable solids.
    \item \textbf{Locality and parallelism:} Interactions have compact support
    and are mostly local, making SPH suitable for parallel CPU and GPU
    implementations.
\end{itemize}

\subsection{Limitations of SPH}

\begin{itemize}
    \item \textbf{Computational cost:} Neighbor searches must be updated as
    particles move, and a well-resolved three-dimensional simulation can
    require many particles and interactions.
    \item \textbf{Accuracy and consistency:} Kernel approximations are
    generally less accurate per degree of freedom than high-order mesh
    methods. Irregular or incomplete particle neighborhoods, especially near
    boundaries, can reduce consistency and introduce error.
    \item \textbf{Pressure and density errors:} Weakly compressible SPH permits
    small density fluctuations and may require very small time steps because
    of an artificial speed of sound. Incompressible SPH variants instead
    require a more expensive global pressure solve.
    \item \textbf{Stability artifacts:} Poor formulations can suffer from
    tensile instability, particle clumping, noisy pressure fields, or
    excessive artificial viscosity. Stabilization improves robustness but can
    add dissipation or alter the modeled physics.
    \item \textbf{Boundary treatment:} Solid walls are not automatically
    represented accurately. Ghost particles, boundary particles, or corrected
    kernels are often needed to impose no-penetration and no-slip conditions.
    \item \textbf{Resolution and convergence:} Refinement, adaptive resolution,
    and rigorous convergence studies are more difficult than in many
    established mesh-based schemes.
\end{itemize}

SPH therefore occupies an important but application-dependent part of the
fluid-simulation field. It is particularly attractive when free surfaces,
moving geometry, violent motion, or fluid fragmentation dominate the problem.
For smooth flows in fixed domains, boundary-layer accuracy, or simulations
requiring high accuracy at lower computational cost, finite volume, finite
element, or high-order finite difference methods are often more appropriate.
Hybrid particle--grid methods are frequently chosen when particle-based
advection is desirable but a grid-based pressure solve or derivative operator
offers better stability.

\section{Context in Fluid Simulation: Functions and Partial Differential Operators}

Fluid motion is governed by continuous continuum mechanics equations, primarily
the \textbf{Navier--Stokes equations} for momentum conservation and the
\textbf{continuity equation} for mass conservation.

\subsection{Governing Equations}

For a standard single-phase Newtonian fluid, the governing model consists of
conservation laws for mass and momentum, constitutive and closure relations,
and appropriate initial and boundary conditions. Thermal flows additionally
require an energy equation. Multiphase, reactive, viscoelastic, or
magnetohydrodynamic simulations require further equations beyond those listed
here.

\subsubsection{Lagrangian Description and Material Derivative}

Because SPH particles move with the fluid, SPH naturally uses the Lagrangian
form of the equations. The material derivative follows a moving fluid parcel:

\begin{equation}
    \frac{D}{Dt}
    = \frac{\partial}{\partial t}
    + \mathbf{v}\cdot\nabla.
\end{equation}

Particle trajectories satisfy the kinematic equation

\begin{equation}
    \frac{D\mathbf{x}}{Dt} = \mathbf{v}.
\end{equation}

\subsubsection{Conservation of Mass}

The continuity equation in Lagrangian form is

\begin{equation}
    \frac{D\rho}{Dt} = -\rho (\nabla \cdot \mathbf{v})
\end{equation}

or, equivalently,

\begin{equation}
    \frac{D\rho}{Dt}+\rho\nabla\cdot\mathbf{v}=0.
\end{equation}

For an incompressible fluid of constant density $\rho_0$, this reduces to the
kinematic constraint

\begin{equation}
    \nabla\cdot\mathbf{v}=0,
    \qquad \rho=\rho_0.
\end{equation}

In SPH, density may be evolved with the continuity equation or recomputed using
the summation density $\rho_i=\sum_j m_jW_{ij}$.

\subsubsection{Conservation of Linear Momentum}

The Cauchy momentum equation is

\begin{equation}
    \rho\frac{D\mathbf{v}}{Dt}
    = \nabla\cdot\boldsymbol{\sigma}
    + \rho\mathbf{b},
\end{equation}

where $\boldsymbol{\sigma}$ is the Cauchy stress tensor and $\mathbf{b}$ is
body force per unit mass, such as gravitational acceleration. Decomposing the
stress into pressure and deviatoric parts gives

\begin{equation}
    \boldsymbol{\sigma}=-p\mathbf{I}+\boldsymbol{\tau}.
\end{equation}

For a Newtonian fluid,

\begin{equation}
    \boldsymbol{\tau}
    = \mu\left(\nabla\mathbf{v}
    +(\nabla\mathbf{v})^{T}\right)
    +\lambda(\nabla\cdot\mathbf{v})\mathbf{I},
\end{equation}

where $\mu$ is dynamic shear viscosity and $\lambda$ is the second viscosity
coefficient. Substitution gives the compressible Navier--Stokes momentum
equation for constant viscosities:

\begin{equation}
    \rho\frac{D\mathbf{v}}{Dt}
    =-\nabla p
    +\mu\nabla^2\mathbf{v}
    +(\mu+\lambda)\nabla(\nabla\cdot\mathbf{v})
    +\rho\mathbf{b}.
\end{equation}

For an incompressible fluid, $\nabla\cdot\mathbf{v}=0$, so this becomes

\begin{equation}
    \rho \frac{D\mathbf{v}}{Dt}
    = -\nabla p + \mu \nabla^2 \mathbf{v}
    + \rho\mathbf{b}+\mathbf{f}_{s},
\end{equation}

where $\mathbf{f}_{s}$ denotes an optional interfacial force per unit volume,
such as surface tension. Angular momentum is conserved when the stress tensor
is symmetric, as it is for an ordinary Newtonian fluid without couple stresses.

\subsubsection{Pressure Closure}

Mass and momentum conservation do not by themselves determine pressure.
The closure depends on the SPH formulation:

\begin{itemize}
    \item \textbf{Weakly compressible SPH (WCSPH):} Pressure is computed from
    an equation of state. A common choice is Tait's equation,
    \begin{equation}
        p = \frac{\rho_0c_0^2}{\gamma}
        \left[
        \left(\frac{\rho}{\rho_0}\right)^\gamma-1
        \right],
    \end{equation}
    where $\rho_0$ is reference density, $c_0$ is an artificial numerical
    speed of sound, and $\gamma$ controls stiffness. The fluid is allowed a
    small, controlled density variation.
    \item \textbf{Incompressible or projection SPH (ISPH):} Pressure acts as a
    Lagrange multiplier enforcing $\nabla\cdot\mathbf{v}=0$. Taking the
    divergence of a provisional-velocity correction produces a pressure
    Poisson equation, schematically
    \begin{equation}
        \nabla\cdot\left(\frac{1}{\rho}\nabla p\right)
        = \frac{1}{\Delta t}\nabla\cdot\mathbf{v}^{*},
    \end{equation}
    followed by
    \begin{equation}
        \mathbf{v}^{n+1}
        =\mathbf{v}^{*}
        -\Delta t\frac{\nabla p}{\rho}.
    \end{equation}
\end{itemize}

\subsubsection{Conservation of Energy}

An isothermal simulation can omit an explicit energy equation. When
temperature, compressible thermodynamics, or heat transfer matters, one common
internal-energy form is

\begin{equation}
    \rho\frac{De}{Dt}
    = -p\nabla\cdot\mathbf{v}
    +\boldsymbol{\tau}:\nabla\mathbf{v}
    -\nabla\cdot\mathbf{q}
    +\rho r,
\end{equation}

where $e$ is specific internal energy, $\mathbf{q}$ is heat flux, and $r$ is
volumetric heating per unit mass. Fourier's law closes conductive heat flux:

\begin{equation}
    \mathbf{q}=-k\nabla T,
\end{equation}

where $k$ is thermal conductivity. A caloric equation of state, such as
$e=e(\rho,T)$, is then also required.

\subsubsection{Initial, Boundary, and Interface Conditions}

The differential equations are not a complete simulation problem until they
are supplied with:

\begin{itemize}
    \item initial particle positions and initial fields
    $\rho(\mathbf{x},0)$, $\mathbf{v}(\mathbf{x},0)$, and, when needed,
    $T(\mathbf{x},0)$;
    \item wall conditions, normally no penetration and either no slip,
    prescribed slip, or a specified wall velocity;
    \item inlet and outlet conditions, such as prescribed velocity, mass flux,
    or pressure;
    \item free-surface conditions, usually ambient normal stress and zero
    tangential stress when surrounding-gas dynamics are neglected; and
    \item interface jump conditions for multiphase flow, including surface
    tension, contact angle, and discontinuous material properties.
\end{itemize}

\subsubsection{Fields and Differential Operators}

In these equations:

\begin{enumerate}
    \item \textbf{Functions ($f(\mathbf{x})$, $g(\mathbf{x})$):} These represent
    continuous fluid properties across space and time, such as the velocity
    field $\mathbf{v}$, density field $\rho$, or pressure field $p$.
    \item \textbf{Partial differential operators
    ($\nabla$, $\nabla \cdot$, $\nabla^2$):} These describe how fluid
    properties change across space.
\end{enumerate}

\begin{itemize}
    \item \textbf{Gradient ($\nabla p$):} Measures spatial pressure
    differences, which generate the internal force that pushes fluid from high
    pressure to low pressure.
    \item \textbf{Divergence ($\nabla \cdot \mathbf{v}$):} Measures the rate at
    which fluid expands or compresses.
    \item \textbf{Laplacian ($\nabla^2 \mathbf{v}$):} Measures spatial
    variations in velocity to compute internal viscous (friction) forces
    between fluid layers.
\end{itemize}

\section{Main Adjustable Parameters in an SPH Simulation}

SPH parameters fall into two fundamentally different groups.
\textbf{Physical parameters} define the fluid and should be determined from the
modeled material and nondimensional regime. \textbf{Numerical parameters}
control how that model is discretized and solved; they should be selected
through convergence, stability, and sensitivity studies rather than tuned to
produce a preferred visual result.

\subsection{Physical and Model Parameters}

\begin{itemize}
    \item \textbf{Reference density $\rho_0$:} Sets the fluid's nominal
    density and, together with particle volume, determines particle mass.
    Water near room temperature is commonly modeled with
    $\rho_0\approx1000\ \mathrm{kg\,m^{-3}}$.
    \item \textbf{Viscosity $\mu$ or $\nu=\mu/\rho$:} Controls physical
    momentum diffusion. Its importance relative to inertia is measured by the
    Reynolds number
    \begin{equation}
        \mathrm{Re}=\frac{UL}{\nu}.
    \end{equation}
    \item \textbf{Body acceleration $\mathbf{b}$:} Usually gravity
    $\mathbf{g}$, but it may also include rotating-frame or prescribed
    accelerations. Gravity--inertia balance is characterized by the Froude
    number $\mathrm{Fr}=U/\sqrt{gL}$.
    \item \textbf{Surface tension $\sigma$:} Controls capillary forces at a
    free surface or material interface. Its importance can be described by the
    Weber number $\mathrm{We}=\rho U^2L/\sigma$.
    \item \textbf{Thermal and multiphase properties:} Thermal conductivity
    $k$, heat capacity, phase densities and viscosities, contact angle, and
    equations of state are adjustable only when the modeled physics requires
    them.
\end{itemize}

Here $U$ and $L$ are characteristic velocity and length scales. Matching the
relevant nondimensional numbers is often more informative than considering
each dimensional parameter independently.

\subsection{Resolution, Particles, and Kernel}

\begin{itemize}
    \item \textbf{Initial particle spacing $\Delta x$:} The primary spatial
    resolution control. In $d$ dimensions, a uniform equal-mass setup commonly
    uses
    \begin{equation}
        m \approx \rho_0(\Delta x)^d.
    \end{equation}
    Halving $\Delta x$ increases particle count by roughly $2^d$, before the
    additional cost of smaller time steps is considered.
    \item \textbf{Smoothing length $h$:} Sets the kernel scale and is usually
    related to particle spacing by
    \begin{equation}
        h=\eta\Delta x,
    \end{equation}
    where $\eta$ is typically of order one but depends on the kernel convention
    and dimension. Increasing $h$ gives more neighbors and smoother, more
    robust estimates, but blurs small features and increases cost. Too small an
    $h$ produces noisy, inconsistent estimates and possible particle disorder.
    \item \textbf{Kernel function and support radius:} Cubic spline, Wendland,
    and other kernels differ in smoothness, stability, pairing resistance, and
    computational cost. A kernel may have support $h$, $2h$, or another
    multiple of $h$, so values of $h$ are meaningful only with the stated
    convention.
    \item \textbf{Target neighbor count:} Closely coupled to $h/\Delta x$ and
    the kernel. Too few neighbors increases sampling error; too many increases
    smoothing and work. Variable-smoothing-length methods adjust $h_i$ to keep
    mass or neighbor count within the support approximately constant.
\end{itemize}

\subsection{Compressibility and Pressure Parameters}

\begin{itemize}
    \item \textbf{Numerical speed of sound $c_0$ in WCSPH:} Determines
    equation-of-state stiffness. A common design criterion is a numerical Mach
    number
    \begin{equation}
        \mathrm{Ma}=\frac{U_{\max}}{c_0}\lesssim0.1,
    \end{equation}
    which keeps density variation near one percent in many flows. A larger
    $c_0$ reduces density error but forces a smaller acoustic time step.
    \item \textbf{Equation-of-state exponent $\gamma$ and background
    pressure:} Control the pressure--density response. For water-like WCSPH,
    $\gamma=7$ is common, but it is a modeling choice rather than a universal
    constant. Background pressure can reduce tensile instability but may
    increase boundary errors or stiffness.
    \item \textbf{Pressure-solver tolerance and iteration limit:} In ISPH,
    PCISPH, DFSPH, and related methods, tighter density or divergence tolerances
    improve incompressibility at greater cost. The iteration cap must be high
    enough that ordinary steps actually reach the tolerance; otherwise it
    silently becomes the effective accuracy control.
\end{itemize}

\subsection{Viscosity, Stabilization, and Particle Regularization}

\begin{itemize}
    \item \textbf{Physical-viscosity discretization:} The chosen Laplacian or
    pairwise formulation must reproduce the intended $\mu$ or $\nu$ and behave
    correctly near walls.
    \item \textbf{Artificial-viscosity coefficients $\alpha$ and $\beta$:}
    Add dissipation to capture shocks and suppress interparticle velocity
    noise. They are important in compressible SPH but can overdamp low-speed
    free-surface flow if set unnecessarily high.
    \item \textbf{Density diffusion, such as $\delta$-SPH:} A diffusion
    coefficient damps high-frequency density and pressure oscillations. Too
    much diffusion smears real density structure and interfaces.
    \item \textbf{XSPH velocity smoothing, shifting, or transport velocity:}
    Regularizes the particle distribution. Its strength trades improved
    particle order and stability against added diffusion or unintended
    transport, so conservation implications must be checked.
\end{itemize}

\subsection{Time Integration and Stability Controls}

\begin{itemize}
    \item \textbf{Time step $\Delta t$:} Must resolve acoustic propagation,
    acceleration, viscosity, and sometimes surface tension. A typical explicit
    choice is
    \begin{equation}
        \Delta t
        = C_{\mathrm{CFL}}\min\left(
        \frac{h}{c_0+v_{\mathrm{sig}}},
        \sqrt{\frac{h}{a_{\max}}},
        \frac{h^2}{\nu}
        \right),
    \end{equation}
    with an additional capillary restriction when surface tension is strong.
    Here $v_{\mathrm{sig}}$ is a characteristic relative or signal velocity
    and $a_{\max}$ is maximum particle acceleration.
    \item \textbf{CFL or safety coefficient $C_{\mathrm{CFL}}$:} Scales the
    stability limits. A smaller value is safer and more accurate but increases
    the number of steps.
    \item \textbf{Integrator:} Explicit Euler, leapfrog, Verlet, predictor--
    corrector, and Runge--Kutta methods have different accuracy, storage, and
    stability properties. Solver tolerances and substepping are additional
    controls for implicit or iterative schemes.
\end{itemize}

\subsection{Boundary and Interaction Parameters}

\begin{itemize}
    \item \textbf{Boundary representation and spacing:} Ghost, dummy, dynamic,
    or semi-analytical boundary methods require enough boundary layers and
    spacing compatible with the fluid particles and kernel support.
    \item \textbf{Wall velocity and slip model:} Determine whether the physical
    condition is no slip, free slip, or partial slip. Penalty-based boundaries
    may also expose stiffness and damping coefficients.
    \item \textbf{Contact and wetting controls:} Contact angle, adhesion, and
    fluid--solid coupling coefficients matter in capillary and multiphysics
    problems and should correspond to measurable physics where possible.
\end{itemize}

\subsection{Practical Parameter-Selection Order}

A defensible setup usually follows this order:

\begin{enumerate}
    \item Set the physical scale, material properties, forcing, and boundary
    conditions; identify the relevant Reynolds, Froude, Weber, and Mach
    numbers.
    \item Choose $\Delta x$ from the smallest feature that must be resolved,
    then choose the kernel and a compatible $h/\Delta x$ and neighbor count.
    \item Select the pressure treatment. For WCSPH, choose $c_0$ from the
    allowable density variation; for an incompressible solver, choose explicit
    density or divergence tolerances.
    \item Add only the stabilization needed for the observed numerical mode,
    and separate artificial dissipation from the target physical viscosity.
    \item Compute $\Delta t$ from all applicable stability restrictions rather
    than selecting it independently.
    \item Repeat the simulation at finer $\Delta x$, smaller $\Delta t$, and
    tighter solver tolerances. Quantities of interest should approach a stable
    value; otherwise the result is parameter-dependent rather than demonstrably
    converged.
\end{enumerate}

\section{Discretization and Its Purpose}

\subsection{What Is Discretization?}

Discretization is the mathematical process of converting \textbf{continuous}
differential equations and infinite spatial fields into a \textbf{discrete},
finite form that a digital computer can compute step by step.

In SPH, discretization happens through a two-step pipeline:

\begin{enumerate}
    \item \textbf{Continuous integral approximation:} Continuous functions are
    smoothed by convolving them with a continuous weighting or smoothing kernel
    function $W$.
    \item \textbf{Particle discretization:} Continuous spatial integrals are
    replaced by discrete summations over nearby neighboring particles:
\end{enumerate}

\subsection{How Kernel Smoothing Produces a Continuous Approximation}

\subsubsection{Field Quantities}

A \textbf{field} assigns a physical quantity to every point in space and,
usually, time. In general, a field can be written as

\begin{equation}
    A = A(\mathbf{x},t),
\end{equation}

where $\mathbf{x}$ is a spatial position and $t$ is time. Holding time fixed
gives the spatial field $A(\mathbf{x})$. For example, the scalar field
$\rho(\mathbf{x},t)$ assigns one density value to each point, while the vector
field $\mathbf{v}(\mathbf{x},t)$ assigns a velocity vector to each point.

\begin{itemize}
    \item \textbf{Scalar fields} assign one number to each location. Examples
    include density $\rho$, pressure $p$, and temperature $T$.
    \item \textbf{Vector fields} assign a vector to each location. Examples
    include velocity $\mathbf{v}$ and force per unit volume $\mathbf{f}$.
\end{itemize}

In continuum mechanics, these fields are treated as continuous: a value is
defined at every point, including the infinitely many points between any two
locations. An SPH computer model cannot store infinitely many values. It stores
field samples $A_j = A(\mathbf{x}_j)$ only at particle positions
$\mathbf{x}_j$ and reconstructs approximate field values between particles.

\subsubsection{The Dirac Delta Distribution}

The \textbf{Dirac delta} $\delta(\mathbf{x}-\mathbf{x}')$ is an idealized point
selector. It is called a distribution rather than an ordinary function because
it is defined by how it behaves inside an integral. Informally, it is zero
everywhere except at $\mathbf{x}'=\mathbf{x}$, is infinitely concentrated at
that point, and has unit integral:

\begin{equation}
    \delta(\mathbf{x}-\mathbf{x}') = 0
    \quad \text{for } \mathbf{x}' \ne \mathbf{x},
    \qquad
    \int_{\Omega} \delta(\mathbf{x}-\mathbf{x}')\,d\mathbf{x}' = 1,
\end{equation}

provided that $\mathbf{x}$ lies inside $\Omega$. Its defining
\textbf{sifting property} is

\begin{equation}
    \int_{\Omega} A(\mathbf{x}')
    \delta(\mathbf{x}-\mathbf{x}')\,d\mathbf{x}' = A(\mathbf{x}).
\end{equation}

Here, $\mathbf{x}$ is the location where the field is being evaluated, while
$\mathbf{x}'$ is a dummy integration variable that ranges over the domain.
The delta eliminates every contribution except the one at
$\mathbf{x}'=\mathbf{x}$, so the integral returns the exact field value at
$\mathbf{x}$. Thus, the exact value of a field can be written as

\begin{equation}
    A(\mathbf{x})
    = \int_{\Omega} A(\mathbf{x}')
    \delta(\mathbf{x} - \mathbf{x}')\,d\mathbf{x}'.
\end{equation}

This identity is the starting point for SPH smoothing. SPH replaces the
infinitely concentrated Dirac delta with a smooth, finite-width kernel
$W(\mathbf{x} - \mathbf{x}', h)$ of finite width. This gives the kernel
approximation

\begin{equation}
    \left\langle A(\mathbf{x}) \right\rangle
    = \int_{\Omega} A(\mathbf{x}')
    W(\mathbf{x} - \mathbf{x}', h)\,d\mathbf{x}'.
\end{equation}

This operation is called a \textbf{convolution}. To estimate the field at
$\mathbf{x}$, it forms a weighted average of values in a neighborhood around
$\mathbf{x}$. Nearby points receive greater weight, while points outside the
kernel's support receive no weight. A useful smoothing kernel normally has the
following properties:

\begin{itemize}
    \item \textbf{Normalization:}
    \begin{equation}
        \int_{\mathbb{R}^d} W(\mathbf{r},h)\,d\mathbf{r} = 1.
    \end{equation}
    This ensures that smoothing reproduces a constant field exactly.
    \item \textbf{Dirac-delta limit:}
    \begin{equation}
        \lim_{h\to 0} W(\mathbf{r},h) = \delta(\mathbf{r}).
    \end{equation}
    More precisely, this limit holds in the distributional sense:
    $\int A(\mathbf{x}')W(\mathbf{x}-\mathbf{x}',h)\,d\mathbf{x}'
    \to A(\mathbf{x})$ for a sufficiently regular field $A$. Thus, the
    approximation converges toward the original field as the smoothing region
    shrinks.
    \item \textbf{Positivity:}
    \begin{equation}
        W(\mathbf{r},h) \geq 0.
    \end{equation}
    Nonnegative weights make the approximation a local weighted average and
    prevent the kernel itself from introducing new extrema through negative
    contributions.
    \item \textbf{Symmetry:}
    \begin{equation}
        W(\mathbf{r},h) = W(-\mathbf{r},h).
    \end{equation}
    A radial kernel depends only on $\lVert\mathbf{r}\rVert$, so particles at
    equal distances receive equal weights regardless of direction. This also
    supports pairwise, antisymmetric force calculations in SPH.
    \item \textbf{Compact support:} There is a finite constant $\kappa$ such
    that
    \begin{equation}
        W(\mathbf{r},h) = 0
        \quad \text{when} \quad
        \lVert\mathbf{r}\rVert \geq \kappa h.
    \end{equation}
    Only nearby particles contribute, making neighbor summations practical.
\end{itemize}

The \textbf{smoothing length} $h$ controls the width of the neighborhood. A
larger $h$ includes more neighbors and produces a smoother but less localized
estimate. A smaller $h$ preserves finer detail but may produce a noisy estimate
if too few particles contribute.

\subsubsection{The Cubic Spline Kernel}

A common SPH choice is the compactly supported cubic spline, also called the
cubic $B$-spline or $M_4$ kernel. Let

\begin{equation}
    q = \frac{\lVert\mathbf{r}\rVert}{h}.
\end{equation}

In $d$ spatial dimensions, the convention used here is

\begin{equation}
    W(\mathbf{r},h) = \sigma_d
    \begin{cases}
        6(q^3-q^2)+1,
            & 0 \leq q \leq \dfrac{1}{2}, \\
        2(1-q)^3,
            & \dfrac{1}{2} < q \leq 1, \\
        0,
            & \text{otherwise},
    \end{cases}
\end{equation}

where the dimension-dependent normalization constant is

\begin{equation}
    \sigma_1 = \frac{4}{3h},
    \qquad
    \sigma_2 = \frac{40}{7\pi h^2},
    \qquad
    \sigma_3 = \frac{8}{\pi h^3}.
\end{equation}

Some references instead define the same cubic spline with support radius
$2h$. The two formulas differ only in the convention used for the smoothing
length; the present form has support radius $h$ and matches the slide notation.

This kernel fulfills the desired conditions:

\begin{itemize}
    \item \textbf{Normalization:} For a radial kernel, changing variables from
    $r$ to $q=r/h$ gives, with $f(q)$ denoting the piecewise polynomial above,
    \begin{equation}
        \int_{\mathbb{R}^d} W\,d\mathbf{r}
        = S_d\sigma_d h^d \int_0^1 f(q)q^{d-1}\,dq,
    \end{equation}
    where $S_1=2$, $S_2=2\pi$, and $S_3=4\pi$ are the surface measures of the
    unit spheres. The needed integrals are
    \begin{equation}
        \int_0^1 f(q)\,dq=\frac{3}{8},
        \qquad
        \int_0^1 qf(q)\,dq=\frac{7}{80},
        \qquad
        \int_0^1 q^2f(q)\,dq=\frac{1}{32}.
    \end{equation}
    Substitution of $\sigma_1$, $\sigma_2$, or $\sigma_3$ therefore makes the
    integral equal to one in the corresponding dimension.
    \item \textbf{Dirac-delta limit:} Each $\sigma_d$ contains the required
    factor $h^{-d}$, preserving unit integral while the support radius $h$
    approaches zero. Consequently, all
    of the kernel's unit mass becomes concentrated at $\mathbf{r}=\mathbf{0}$,
    $W(\mathbf{r},h) \rightarrow \delta(\mathbf{r})$ as $h \rightarrow 0$, so
    convolution with $W$ approaches exact evaluation.
    \item \textbf{Positivity:} Both polynomial pieces of $f(q)$ are
    nonnegative on their stated intervals, and $\sigma_d>0$. Hence
    $W(\mathbf{r},h)\geq 0$.
    \item \textbf{Symmetry:} The kernel depends on $\mathbf{r}$ only through
    $q=\lVert\mathbf{r}\rVert/h$. Since
    $\lVert-\mathbf{r}\rVert=\lVert\mathbf{r}\rVert$, it follows that
    $W(-\mathbf{r},h)=W(\mathbf{r},h)$.
    \item \textbf{Compact support:} The kernel is zero outside $q\leq1$ and
    also reaches zero at $q=1$. Therefore $W(\mathbf{r},h)=0$ whenever
    $\lVert\mathbf{r}\rVert\geq h$, so $\kappa=1$.
\end{itemize}

To evaluate the integral using particles, SPH assigns particle $j$ the volume
element

\begin{equation}
    \Delta V_j \approx \frac{m_j}{\rho_j},
\end{equation}

where $m_j$ and $\rho_j$ are its mass and density. Replacing the integral by a
sum of these particle volumes gives

\begin{equation}
    \left\langle A(\mathbf{x}_i) \right\rangle
    \approx \sum_j A_j \frac{m_j}{\rho_j}
    W(\mathbf{x}_i - \mathbf{x}_j, h)
\end{equation}

Each term is particle $j$'s field value multiplied by its represented volume
and its kernel weight relative to particle $i$. Although the field data is
stored only at discrete particles, $W$ is a continuous function of position.
The same sum can therefore be evaluated at any location $\mathbf{x}$, producing
a continuous interpolated field from discrete samples.

Similarly, spatial derivatives on fluid particles shift directly onto
gradients of the kernel function $W$.

\subsection{Worked PDE Example: The Continuity Equation}

Consider conservation of mass in Eulerian form:

\begin{equation}
    \frac{\partial \rho}{\partial t}
    + \nabla\cdot(\rho\mathbf{v}) = 0.
\end{equation}

Using the material derivative

\begin{equation}
    \frac{D}{Dt}
    = \frac{\partial}{\partial t} + \mathbf{v}\cdot\nabla,
\end{equation}

the same PDE can be written in Lagrangian form as

\begin{equation}
    \frac{D\rho}{Dt} = -\rho\nabla\cdot\mathbf{v}.
\end{equation}

This form is natural for SPH because each particle moves with the fluid.

\subsubsection{How SPH Processes the Field Quantities}

The continuous fields are replaced by values carried by particles:

\begin{itemize}
    \item The density field $\rho(\mathbf{x},t)$ becomes one scalar sample
    $\rho_i(t)$ on particle $i$.
    \item The velocity field $\mathbf{v}(\mathbf{x},t)$ becomes one vector
    sample $\mathbf{v}_i(t)$ on particle $i$.
    \item The particle position $\mathbf{x}_i(t)$ specifies where those field
    samples are currently located.
    \item The mass $m_i$ is usually fixed and represents the amount of fluid
    carried by the particle. The ratio $m_i/\rho_i$ approximates its volume.
\end{itemize}

When a field value is required between particles, SPH reconstructs it using
the kernel sum. For example,

\begin{equation}
    \left\langle \rho(\mathbf{x}) \right\rangle
    \approx \sum_j m_j W(\mathbf{x}-\mathbf{x}_j,h).
\end{equation}

At particle $i$, this gives the commonly used density summation

\begin{equation}
    \rho_i \approx \sum_j m_j W_{ij},
    \qquad
    W_{ij}=W(\mathbf{x}_i-\mathbf{x}_j,h).
\end{equation}

Thus, field processing means storing samples on particles and using weighted
neighbor sums whenever a continuous value must be reconstructed.

\subsubsection{How SPH Processes the Differential Operators}

SPH replaces each differential operator with an operation on moving particles
and kernel derivatives:

\begin{itemize}
    \item The material derivative $D/Dt$ becomes an ordinary time derivative
    following particle $i$. Therefore
    $D\rho/Dt$ becomes $d\rho_i/dt$.
    \item The divergence $\nabla\cdot\mathbf{v}$ becomes a weighted sum of
    velocity differences dotted with kernel gradients:
\end{itemize}

\begin{equation}
    (\nabla\cdot\mathbf{v})_i
    \approx
    \frac{1}{\rho_i}
    \sum_j m_j(\mathbf{v}_j-\mathbf{v}_i)
    \cdot\nabla_i W_{ij}.
\end{equation}

The difference $\mathbf{v}_j-\mathbf{v}_i$ is important: if all neighboring
particles have the same velocity, every difference is zero and SPH correctly
estimates zero divergence.

For the cubic spline used above, let
$\mathbf{r}_{ij}=\mathbf{x}_i-\mathbf{x}_j$,
$r_{ij}=\lVert\mathbf{r}_{ij}\rVert$, and $q=r_{ij}/h$. For
$r_{ij}>0$, its gradient is

\begin{equation}
    \nabla_i W_{ij}
    = \frac{\sigma_d}{h}f'(q)
    \frac{\mathbf{r}_{ij}}{r_{ij}},
\end{equation}

where

\begin{equation}
    f'(q)=
    \begin{cases}
        18q^2-12q,
            & 0 \leq q \leq \dfrac{1}{2}, \\
        -6(1-q)^2,
            & \dfrac{1}{2}<q\leq1, \\
        0,
            & \text{otherwise}.
    \end{cases}
\end{equation}

At $r_{ij}=0$, symmetry gives $\nabla_iW_{ij}=\mathbf{0}$.

\subsubsection{The Resulting Particle Equation}

Substituting the SPH divergence approximation into the continuity equation
gives

\begin{align}
    \frac{d\rho_i}{dt}
    &= -\rho_i(\nabla\cdot\mathbf{v})_i \\
    &\approx \sum_j m_j(\mathbf{v}_i-\mathbf{v}_j)
    \cdot\nabla_iW_{ij}.
\end{align}

The continuous PDE has now become one ordinary differential equation for each
particle. A simulation step processes it as follows:

\begin{enumerate}
    \item Find all neighbors $j$ satisfying
    $\lVert\mathbf{x}_i-\mathbf{x}_j\rVert<h$.
    \item Read the sampled field values $\rho_i$, $\mathbf{v}_i$, and
    $\mathbf{v}_j$ from the particles.
    \item Evaluate $\nabla_iW_{ij}$ from particle separation using the cubic
    spline derivative.
    \item Sum the neighbor contributions to obtain $d\rho_i/dt$.
    \item Advance density in time; for example, an explicit Euler step is
    \begin{equation}
        \rho_i^{n+1}
        = \rho_i^n + \Delta t
        \left(\frac{d\rho_i}{dt}\right)^n.
    \end{equation}
\end{enumerate}

In summary, SPH handles \emph{field quantities} as particle-carried samples and
kernel-weighted reconstructions. It handles \emph{differential operators} by
differentiating the smooth kernel and summing the resulting neighbor
contributions. This converts the original spatial PDE into a system of
time-dependent particle equations.

\section{Major SPH Papers and Further Reading}

The following papers provide a chronological path from the original
astrophysical formulation of SPH to modern incompressible-flow and computer
graphics methods. DOI links lead to the publisher record; access to the full
text may depend on the reader's institution.

\begin{enumerate}
    \item \textbf{L. B. Lucy (1977), ``A Numerical Approach to the Testing of
    the Fission Hypothesis.''} One of the two papers that independently
    introduced SPH, originally for three-dimensional astrophysical problems.
    \href{https://doi.org/10.1086/112164}{Read via DOI}.

    \item \textbf{R. A. Gingold and J. J. Monaghan (1977), ``Smoothed Particle
    Hydrodynamics: Theory and Application to Non-Spherical Stars.''} The other
    foundational SPH paper; it develops the kernel-based particle formulation
    and applies it to stellar dynamics.
    \href{https://doi.org/10.1093/mnras/181.3.375}{Read via DOI}.

    \item \textbf{J. J. Monaghan (1992), ``Smoothed Particle Hydrodynamics.''}
    A classic review that systematizes the mathematical formulation, kernels,
    conservation properties, artificial viscosity, and early astrophysical
    applications.
    \href{https://doi.org/10.1146/annurev.aa.30.090192.002551}{Read via DOI}.

    \item \textbf{J. P. Morris, P. J. Fox, and Y. Zhu (1997), ``Modeling Low
    Reynolds Number Incompressible Flows Using SPH.''} An influential treatment
    of physical viscosity and wall boundaries for viscous incompressible flows.
    \href{https://doi.org/10.1006/jcph.1997.5776}{Read via DOI}.

    \item \textbf{S. J. Cummins and M. Rudman (1999), ``An SPH Projection
    Method.''} Introduces a pressure-projection approach that enforces
    incompressibility through a pressure Poisson equation rather than a stiff
    equation of state.
    \href{https://doi.org/10.1006/jcph.1999.6246}{Read via DOI}.

    \item \textbf{M. M{\"u}ller, D. Charypar, and M. Gross (2003),
    ``Particle-Based Fluid Simulation for Interactive Applications.''} A
    landmark computer-graphics paper that presents an efficient SPH
    formulation for interactive liquid simulation, including pressure,
    viscosity, and surface tension.
    \href{https://doi.org/10.2312/SCA03/154-159}{Read via DOI}.

    \item \textbf{J. J. Monaghan (2005), ``Smoothed Particle Hydrodynamics.''}
    A broad and more modern review of SPH theory, numerical behavior, and
    applications across astrophysics and continuum mechanics.
    \href{https://doi.org/10.1088/0034-4885/68/8/R01}{Read via DOI}.

    \item \textbf{M. Becker and M. Teschner (2007), ``Weakly Compressible SPH
    for Free Surface Flows.''} A widely used graphics formulation that controls
    density variation with a stiff equation of state and is designed for
    efficient free-surface animation.
    \href{https://doi.org/10.2312/SCA/SCA07/209-217}{Read via DOI}.

    \item \textbf{B. Solenthaler and R. Pajarola (2009), ``Predictive-Corrective
    Incompressible SPH.''} Introduces PCISPH, which iteratively predicts and
    corrects pressure to reduce density error while allowing larger time steps
    than many weakly compressible formulations.
    \href{https://doi.org/10.1145/1576246.1531346}{Read via DOI}.

    \item \textbf{D. J. Price (2012), ``Smoothed Particle Hydrodynamics and
    Magnetohydrodynamics.''} A detailed derivation-focused review covering
    consistency, conservation, shocks, dissipation, and magnetohydrodynamics;
    it is useful even when the main interest is ordinary fluid SPH.
    \href{https://doi.org/10.1016/j.jcp.2010.12.011}{Read via DOI}.

    \item \textbf{A. Adami, X. Y. Hu, and N. A. Adams (2012), ``A Generalized
    Wall Boundary Condition for Smoothed Particle Hydrodynamics.''} Develops a
    robust boundary-particle treatment for walls under complex and
    time-dependent flows.
    \href{https://doi.org/10.1016/j.jcp.2012.05.005}{Read via DOI}.

    \item \textbf{M. Ihmsen, J. Orthmann, B. Solenthaler, A. Kolb, and
    M. Teschner (2014), ``SPH Fluids in Computer Graphics.''} A useful survey
    of the complete SPH simulation pipeline, including neighborhood search,
    incompressibility solvers, boundary handling, surface reconstruction, and
    implementation considerations.
    \href{https://doi.org/10.1109/TVCG.2014.2309246}{Read via DOI}.
\end{enumerate}

\end{document}